\scanpage{35}
\mspage{\handwritten{- 42 -}}
% archive sheet number: 34

convenable $s\handwritten{\to}s+a$ et la déformation $(t,s)\handwritten{\to\varphi}_2(t,s+a)$ de $\handwritten{\gamma}_1$ en $\handwritten{\gamma}_2$ ; alors si $\handwritten{\varphi}$ est la fonction égale à $\handwritten{\varphi}_1$ dans $I\handwritten{\times}J_1$ et à $\handwritten{\varphi}_2$ dans $I\handwritten{\times}J_2$ (ce qui a un sens; puisque les deux fonctions coïncident sur $I\handwritten{\times}(J_1\handwritten{\cap}J_2)$), il est clair que $\handwritten{\varphi}$ est continue dans $I\handwritten{\times}(J_1\handwritten{\cup}J_2)$ et est une déformation de $\handwritten{\gamma}_0$ en $\handwritten{\gamma}_2$ dans A. Comme en outre tout chemin est déformable en lui-même (en prenant $\handwritten{\varphi}(t,s)=\handwritten{\gamma}(t)$ pour tout $s\handwritten{\in}J$), la relation ``$\handwritten{\gamma}$ est déformable en $\handwritten{\gamma}'$ dans A'' est une \uline{relation d'équivalence} entre les chemins définis dans I et à valeurs dans A.

\uline{Exemple}.- On dit qu'un domaine $A\handwritten{\subset\mathbb{R}}^n$ est \uline{étoilé} par rapport à un point $\handwritten{a\in}A$ si pour tout $\handwritten{x\in}A$, le segment fermé d'extrémités $\handwritten{a}$ et $\handwritten{x}$ est contenu dans A. Dans un domaine étoilé par rapport à A, \uline{tout chemin} (resp. \uline{tout lacet}) $\handwritten{\gamma}$ \uline{est déformable en le point} $\handwritten{a}$ par la déformation $\handwritten{\varphi}(t,s)=\handwritten{a}+(1-s)(\handwritten{\gamma}(t)-\handwritten{a})$ (pour $0\handwritten{\leq}s\handwritten{\leq}1$). Deux chemins (resp. lacets) quelconques sont donc déformables l'un en l'autre dans A.

\uline{Remarque}.- Si $\handwritten{\gamma}$ est une juxta-position de deux chemins $\handwritten{\gamma}_1$, $\handwritten{\gamma}_2$, et $\handwritten{\varphi}$ une déformation de $\handwritten{\gamma}$ en $\handwritten{\gamma}'$ dans A, $\handwritten{\gamma}'$ est juxtaposition des deux chemins $\handwritten{\gamma}'_1$, $\handwritten{\gamma}'_2$ déformés de $\handwritten{\gamma}_1$, $\handwritten{\gamma}_2$ par $\handwritten{\varphi}$. Mais on aura soin d'observer que si $\handwritten{\gamma}_1$ et $\handwritten{\gamma}_2$ peuvent être \uline{chacun} déformés \margin{\handwritten{2}}dans A en deux nouveaux chemins $\handwritten{\gamma}'_1$\unsure{,|} $\handwritten{\gamma}'_2$, tels que l'extrémité de $\handwritten{\gamma}'_1$ soit l'origine de $\handwritten{\gamma}'_2$, il ne s'ensuit nullement que $\handwritten{\gamma}$ puisse être déformé dans A en la juxtaposition de $\handwritten{\gamma}'_1$ et de $\handwritten{\gamma}'_2$.

PROPOSITION 1.- \uline{Tout chemin achevé dans} A \uline{peut être déformé dans} A \uline{en une route linéaire par morceaux ayant mêmes extrémités}.

En effet, soit $\handwritten{\gamma}$ un chemin a chevé à valeurs dans A, défini dans l'intervalle compact $I=\handwritten{[}a,b\handwritten{]}$. L'ensemble $\handwritten{\gamma}(I)$ est compact et contenu dans A ; la distance $\handwritten{\rho}$ de cet ensemble à $\handwritten{\complement}A$ est donc $\handwritten{>}0$.
